\documentclass[12pt]{article}
\usepackage[utf8]{inputenc}
\usepackage[T1]{fontenc}
\usepackage{hyperref}
\title{Kubachi Language Resources for K-12 Curriculum Development}
\author{Compiled by Claude Sonnet 5 (claude-sonnet-5) for Andr\'as Kornai}
\date{}
\begin{document}
\maketitle
\section{Introduction}
Kubachi (also Ughbug) is a Dargwa-group language of Dagestan, Russia. It is often treated as a variant of Dargwa. The search found academic grammar and phonetics work, but no school materials and no Kubachi-specific dictionary.
\section{Primary Educational Resources}
None found.
\section{Linguistic and Reference Materials}
\begin{itemize}
\item Yamling and Tchantouria (1991), a Lund working paper sketching Kubachi simple-sentence grammar. It draws mainly on Magometov's 1963 monograph: \url{https://journals.lub.lu.se/LWPL/article/download/2580/2155/5234}
\item Comparative phonetic study of the Kubachi and Itsari lects: \url{https://gramota.net/eng/materials/2/2023/12/41.html}
\item Research article on Kubachi phonetics (full text): \url{https://philology-journal.ru/en/article/phil20230665/fulltext}
\item BnF catalogue entry listing Ughbug as an alternate name: \url{https://catalogue.bnf.fr/ark:/12148/cb16505644b}
\end{itemize}
Magometov's 1963 monograph was cited in these sources, but no URL for it appeared in the results.
\section{Religious and Community Resources}
None found.
\section{Digital and Audio Resources}
The electronic Dargwa dictionary by Mutalov and Mutalov covers Dargwa generally, not Kubachi specifically. No URL for it appeared in the results.
\section{Curriculum Development Context}
Documentation is mainly academic and Russian-language. No school programme or Kubachi-specific teaching material was found. A curriculum would need to build on the grammars above and on work with the community.
\end{document}